\begin{table}[t]
\centering
\caption{Available water capacity (AWC) for the 10 MATOPIBA municipalities used in the validation grid, derived from the national CAD dataset of the Brazilian Water Agency \citep{SantosANAUFPR_TED}, available at SNIRH \citep{SNIRH28fe4baa}. The CAD volumetric value (m\textsuperscript{3}\,m\textsuperscript{-3}) is the per-municipality average over the 0--120\,cm reference profile, computed by the UFPR/ANA team via Saxton \& Rawls \citep{SaxtonRawls2006} pedotransfer functions calibrated on HYBRAS \citep{Ottoni2018HYBRAS} laboratory profiles and extrapolated through the 1:250{,}000 RADAM/IBGE soil map. The AWC in mm is integrated over the early-cycle soybean root zone \(Z_r=60\)\,cm: AWC\textsubscript{mm} = CAD \(\times Z_r \times 10\). Cities sorted by ascending AWC.}
\label{tab:soils-matopiba}
\begin{tabular}{lrrr}
\toprule
Municipality & IBGE code & CAD (m\textsuperscript{3}\,m\textsuperscript{-3}) & AWC (mm) \\
\midrule
Bom Jesus & 2201903 & 0.0573 & 34.4 \\
Campos Lindos & 1703842 & 0.0608 & 36.5 \\
Urucui & 2211209 & 0.0689 & 41.4 \\
Baixa Grande do Ribeiro & 2201150 & 0.0772 & 46.4 \\
Correntina & 2909307 & 0.0777 & 46.6 \\
Barreiras & 2903201 & 0.0814 & 48.9 \\
Balsas & 2101400 & 0.0837 & 50.2 \\
Luis Eduardo Magalhaes & 2919553 & 0.0881 & 52.9 \\
Tasso Fragoso & 2112001 & 0.0888 & 53.3 \\
Formosa do Rio Preto & 2911105 & 0.1021 & 61.2 \\
\midrule
\textbf{Mean} & & \textbf{0.0786} & \textbf{47.2} \\
\bottomrule
\end{tabular}
\end{table}